%% =====================================================================
%%  T-24-06  Judge the Deed
%%  -------------------------------------------------------------------
%%  One idea per file. Core is mandatory. Slots in canonical order.
%%  Max 700 words. No layout commands. No filler. New source material
%%  for this entry goes in sources/B6/T-24-06.tex -- never by
%%  rewriting this file (docs/RULES.md R0.1).
%%  =====================================================================
%% @kind: topic
%% @id: T-24-06
%% @title: Judge the Deed
%% @subtitle: Intentions are unreadable and irrelevant --- only behaviour is under negotiation
%% @chapter: c24
%% @order: 40
%% @status: stable
%% @sources: B6
%% @updated: 2026-09-19

\topic{T-24-06}{Judge the Deed}{Intentions are unreadable and irrelevant --- only behaviour is under negotiation}

\begin{core}
B6's counter-doctrine, distilled: never mind-read --- you cannot know why someone hurt you, and in the end it is irrelevant. The whole engagement is renegotiated on the only medium that can be observed: behaviour. Accept no excuses, insist on direct answers, discount promises until behaviour changes, and move early --- because the aggressor sets the rules of any fight they are allowed to start.
\end{core}
\begin{mechanism}
The doctrine's first principle is engagement-shape: whoever opens the aggression writes the fight's rules --- once attacked, weakened or put on the run, the victim spends the whole contest scrambling for balance. Hence every tool aims to stop reacting and start pricing. Judge actions, not intentions: the motive is undiscoverable and serves only as a stage for argument. Accept no excuses: an excuse is a prediction --- a person already resisting submission to the principle of civil conduct will repeat the behaviour; drop the excuse and the conduct can be labelled plainly. Accept only direct responses: manipulators are notorious for never giving a straight answer to a straight question, so a direct request earns a request for a direct answer, repeated without hostility, which matters because any sarcasm or put-down is construed as an attack and licenses their war. Discount promises: promises mean nothing until the usual manipulations are visibly interrupted --- changed behaviour is the only evidence of change. And move early: a runaway train stops easiest before it gathers speed; the aggressive personality lacks internal brakes, so the cheap moment is the first moment. The supporting frame: agreements worth making are proportionate, dependable, checkable and backed by consequence --- and stop fighting battles you cannot win, which is where the targets of the world exhaust themselves into depression.
\end{mechanism}
\begin{conditions}
Strongest against the habitually manipulative, where motive-talk is a trap and records are the only level ground. It transfers imperfectly to ordinary relationships: run at full strength on family and friends who are merely clumsy, it reads as cold litigation --- the doctrine is calibrated for patterns, not incidents. It requires the target's own house in order: agreements you will not keep license the other side's exemptions.
\end{conditions}
\begin{application}
  \item Convert every dispute to the deed: what did they do, what did you ask for, what happened next. Motives come off the table by rule, not by hope.
  \item When an excuse arrives, label the conduct once --- \emph{that is what happened} --- and stop litigating the why.
  \item Re-ask direct questions without hostility until answered. The third calm repetition is usually the one that gets an answer.
  \item Price change by behaviour over time, never by the promise made in the repair phase.
  \item Move at the first roll of the train. Every tactic is cheaper to decline in week one than in year five.
\end{application}
\begin{feedback}
  \item Working: your confrontations are shorter and end in agreements rather than explanations --- the motive-stage has been dismantled.
  \item Working: excuses now read as forecasts, and the forecast verifies. The pricing is sound.
  \item Failing: you are running motive-trials again because the deed-record is inconvenient. The doctrine never fails from the deed side.
  \item Failing: everyone experiences you as a contract. The calibration has drifted from patterns to people; reread the conditions.
\end{feedback}
\begin{failure}
The doctrine's cost is warmth at the margin: relationships priced purely in deeds lose the benefit of the doubt that intimacy runs on, and the person who applies litigation to love will win every case and lose every client. It also has a shadow use --- the operator's version, \emph{ignore their stated reasons, just price their behaviour}, is itself how the coldest operators treat everyone; the entry's defence is the conditions: patterns, not incidents; records, not contempt.
\end{failure}
\begin{countermeasures}
  \item Notice who keeps the engagement on your motives --- always your tone, your intent, your sensitivity --- and never on deeds. The venue is chosen; move the venue back.
  \item Write your own agreements down if you expect to honour them. The enforceable clause begins at home (T-04-02).
  \item Distinguish one bad deed from a pattern before deploying the doctrine; otherwise it is just a new way to be cruel efficiently.
  \item When you hear yourself demanding explanations of someone's motives, ask who taught you that the motive mattered more than the act.
\end{countermeasures}

\sources{B6}
\seesources
\seealso{T-24-01, T-06-04, T-23-04}
